\renewcommand{\Anno}{2004}

\renewcommand{\NumeroProblema}{4}

\renewcommand{\Testo}{Sono assegnati tre punti distinti su una
circonferenza.  Ad ogni mossa \`{e} possibile spostare uno qualsiasi
dei tre punti, facendogli occupare come nuova posizione il punto medio
di uno, a scelta, dei due archi di circonferenza determinati dagli
altri due punti.
	
\begin{enumerate}
	\item Determinare per quali configurazioni iniziali esiste una
	successione finita di mosse che li porta alla fine ad essere i
	vertici di un triangolo isoscele con due angoli minori di
	1/2004-esimo di grado.
	
	\item Determinare per quali configurazioni iniziali esiste una
	successione finita di mosse che li porta alla fine ad essere i
	vertici di un triangolo equilatero.
\end{enumerate}}

\renewcommand{\Risposta}{\begin{enumerate}
		\item  Per ogni configurazione iniziale.
	
		\item  Per tutte e sole le configuarazioni iniziali in cui almeno 
		due punti determinano un  angolo al centro di $120^{\circ}$.
	\end{enumerate}}

\renewcommand{\Aiutino}{\begin{enumerate}
	\renewcommand{\labelenumi}{(\alph{enumi})} 
	\item Far vedere che, agendo opportunamente, \`{e} possibile 
	costruire triangoli isosceli con l'angolo alla base che si dimezza di 
	volta in volta.
	
	\item Scrivere nei due casi gli angoli del triangolo finale in 
	funzione degli angoli del triangolo iniziale. Se uno degli angoli nel 
	finale misura $60^{\circ}$, come potevano/dovevano essere gli angoli 
	nel triangolo iniziale?
\end{enumerate}}

\renewcommand{\SoluzioneUfficiale}{\TitoloSoluzione \DomandaA Dopo la
prima mossa i 3 punti sono i vertici di un triangolo isoscele.  Se ora
spostiamo un estremo della base nel punto medio del pi\`{u} piccolo
arco delimitato dagli altri 2 punti, otteniamo un nuovo triangolo
isoscele i cui angoli alla base hanno un'ampiezza che \`{e} la
met\`{a} rispetto agli angoli alla base del triangolo precedente (la 
dimostrazione di questo fatto \`{e} contenuta nella figura centrale 
sottostante, in cui i 4 angoli evidenziati hanno uguale ampiezza 
perch\'{e} insistono su archi congruenti).
\begin{center}
	\hfill
	\psset{unit=8ex}
	\pspicture(-1.2,-1.2)(1.2,1.3)
	
	\pscircle[linecolor=green](0,0){1}
	\pstGeonode[PointSymbol={none,*},PosAngle={0,90,-45,-135}](0,0){O}(1;90){A}(1;-65){C}(1;-115){B}
	\pstRotation[PointSymbol=none,AngleCoef=0.5,RotAngle=\pstAngleAOB{A}{O}{B}]{O}{A}{C'}
	
	\pstLineAB{A}{B}
	\pstLineAB{A}{C}
	\pstLineAB{C}{B}
	\endpspicture
	\hfill
	\pspicture(-1.2,-1.2)(1.2,1.3)
	
	\pscircle[linecolor=green](0,0){1}
	\pstGeonode[PointSymbol={none,*},PosAngle={0,90,-45,-135}](0,0){O}(1;90){A}(1;-65){C}(1;-115){B}
	\pstRotation[PosAngle=160,AngleCoef=0.5,RotAngle=\pstAngleAOB{A}{O}{B}]{O}{A}{C'}
	
	\pstLineAB{A}{B}
	\pstLineAB{A}{C}
	\pstLineAB{C}{B}
	
	\pstLineAB{A}{C'}
	\pstLineAB{B}{C'}
	\pstLineAB{C}{C'}
	
	\psset{MarkAngleRadius=0.25,linecolor=blue}
	\pstMarkAngle{A}{B}{C'}{}
	\pstMarkAngle{C'}{A}{B}{}
	\pstMarkAngle[MarkAngleRadius=0.2]{A}{C}{C'}{}
	\pstMarkAngle{C'}{C}{B}{}
	\endpspicture
	\hfill
	\pspicture(-1.2,-1.2)(1.2,1.3)
	
	\pscircle[linecolor=green](0,0){1}
	\pstGeonode[PointSymbol={none,*,none,*},PosAngle={0,90,-45,-135}](0,0){O}(1;90){A}(1;-65){C}(1;-115){B}
	\pstRotation[PosAngle=160,AngleCoef=0.5,RotAngle=\pstAngleAOB{A}{O}{B}]{O}{A}{C'}
	
	\pstLineAB{A}{B}
	% \pstLineAB{A}{C}
	% \pstLineAB{C}{B}
	
	\pstLineAB{A}{C'}
	\pstLineAB{B}{C'}
	% \pstLineAB{C}{C'}
	
	\psset{MarkAngleRadius=0.25,linecolor=blue}
	\pstMarkAngle{A}{B}{C'}{}
	\pstMarkAngle{C'}{A}{B}{}
	% \pstMarkAngle[MarkAngleRadius=0.2]{A}{C}{C'}{}
	% \pstMarkAngle{C'}{C}{B}{}
	\endpspicture
	\hfill\mbox{}
\end{center}

Continuando ad operare allo stesso modo, ad ogni passaggio l'ampiezza
degli angoli alla base si dimezza.  Possiamo quindi, partendo da
qualunque configurazione, ottenere triangoli isosceli con angoli alla 
base piccoli a piacere. 
\begin{center}
	\hfill
	\psset{unit=8ex}
	\pspicture(-1,-1)(1,1)
	
	\pscircle[linecolor=green](0,0){1}
	\pstGeonode[PointSymbol={none,*},PointName=none](0,0){O}(1;90){A}(1;-65){C}(1;-115){B}
	\pstRotation[PointSymbol=square*,linecolor=blue,PointName=none,AngleCoef=0.5,RotAngle=\pstAngleAOB{A}{O}{B}]{O}{A}{C'}
	
	\pstLineAB{A}{B}
	\pstLineAB{A}{C}
	\pstLineAB{C}{B}
	\endpspicture
	\hfill
	\pspicture(-1,-1)(1,1)
	
	\pscircle[linecolor=green](0,0){1}
	\pstGeonode[PointSymbol={none,*,none,*},PointName=none](0,0){O}(1;90){A}(1;-65){C}(1;-115){B}
	\pstRotation[PointName=none,AngleCoef=0.5,RotAngle=\pstAngleAOB{A}{O}{B}]{O}{A}{C'}
	\pstRotation[PointSymbol=square*,linecolor=blue,PointName=none,AngleCoef=0.5,RotAngle=\pstAngleAOB{A}{O}{C'}]{O}{A}{B'}
	
	\pstLineAB{A}{B}
	\pstLineAB{A}{C'}
	\pstLineAB{C'}{B}
	\endpspicture
	\hfill
	\pspicture(-1,-1)(1,1)
	
	\pscircle[linecolor=green](0,0){1}
	\pstGeonode[PointSymbol={none,*,none},PointName=none](0,0){O}(1;90){A}(1;-65){C}(1;-115){B}
	\pstRotation[PointName=none,AngleCoef=0.5,RotAngle=\pstAngleAOB{A}{O}{B}]{O}{A}{C'}
	\pstRotation[PointName=none,AngleCoef=0.5,RotAngle=\pstAngleAOB{A}{O}{C'}]{O}{A}{B'}
	\pstRotation[PointSymbol=square*,linecolor=blue,PointName=none,AngleCoef=0.5,RotAngle=\pstAngleAOB{A}{O}{B'}]{O}{A}{C''}
	
	\pstLineAB{A}{B'}
	\pstLineAB{A}{C'}
	\pstLineAB{C'}{B'}
	\endpspicture
	\hfill
	\pspicture(-1,-1)(1,1)
	
	\pscircle[linecolor=green](0,0){1}
	\pstGeonode[PointSymbol={none,*,none},PointName=none](0,0){O}(1;90){A}(1;-65){C}(1;-115){B}
	\pstRotation[PointSymbol=none,AngleCoef=0.5,RotAngle=\pstAngleAOB{A}{O}{B}]{O}{A}{C'}
	\pstRotation[PointName=none,AngleCoef=0.5,RotAngle=\pstAngleAOB{A}{O}{C'}]{O}{A}{B'}
	\pstRotation[PointName=none,AngleCoef=0.5,RotAngle=\pstAngleAOB{A}{O}{B'}]{O}{A}{C''}
	
	\pstLineAB{A}{B'}
	\pstLineAB{A}{C''}
	\pstLineAB{C''}{B'}
	\endpspicture
	\hfill
	\mbox{}
\end{center}

\DomandaB Supponiamo di partire da una configurazione data $A$, $B$, 
$C$, e di spostare il punto $C$ nelle 2 posizioni $C'$ possibili e 
determiniamo nei 2 casi gli angoli del triangolo $ABC'$ in funzione 
degli angoli di $ABC$.


Nel primo caso $C$ e $C'$ stanno dalla stessa parte rispetto ad 
$AB$, pertanto $\gamma'=\gamma$ e
$\alpha'=\beta'=(180^{\circ}-\gamma)/2=90^{\circ}-\gamma/2$. Quindi 
in questo caso gli angoli cambiano nel seguente modo
$$(\alpha,\beta,\gamma)\longrightarrow
\left(90^{\circ}-\frac{\gamma}{2},90^{\circ}-\frac{\gamma}{2},
\gamma\right).$$

Nel secondo caso $C$ e $C'$ stanno da parti opposte rispetto ad 
$AB$, pertanto $\gamma'=180^{\circ}-\gamma$, da cui 
$\alpha'=\beta'=\gamma/2$. Quindi gli angoli cambiano nel seguente 
modo
$$(\alpha,\beta,\gamma)\longrightarrow
\left(\frac{\gamma}{2},\frac{\gamma}{2},180^{\circ}-\gamma\right).$$

\begin{center}
	\psset{unit=9ex}
	\pspicture(-1,-3)(6,3)
	
	\rput(0,0){\pspicture(-1,-1)(1,1)
	
	\pscircle[linecolor=green](0,0){1}
	
	\pstGeonode[PosAngle={200,-20}](1;-160){A}(1;-20){B}
	\pstGeonode[PosAngle=130](1;130){C}
	
	\pstLineAB{A}{B}
	\pstLineAB{A}{C}
	\pstLineAB{C}{B}
	
	\pstMarkAngle[MarkAngleRadius=0,LabelSep=0.3]{B}{A}{C}{$\alpha$}
	\pstMarkAngle[MarkAngleRadius=0,LabelSep=0.6]{C}{B}{A}{$\beta$}
	\pstMarkAngle[MarkAngleRadius=0,LabelSep=0.3]{A}{C}{B}{$\gamma$}
	\endpspicture
	}
	
	\rput(5,1.5){\pspicture(-1,-1)(1,1)
	\pscircle[linecolor=green](0,0){1}
	\pstGeonode[PosAngle={200,-20}](1;-160){A}(1;-20){B}
	\pstGeonode[PosAngle=90,PointName=C'](1;90){C}
	
	\pstLineAB{A}{B}
	\pstLineAB{A}{C}
	\pstLineAB{C}{B}
	
	\pstMarkAngle[MarkAngleRadius=0,LabelSep=0.4]{B}{A}{C}{$\alpha'$}
	\pstMarkAngle[MarkAngleRadius=0,LabelSep=0.4]{C}{B}{A}{$\beta'$}
	\pstMarkAngle[MarkAngleRadius=0,LabelSep=0.4]{A}{C}{B}{$\gamma'$}
	\endpspicture
	}
	\rput(5,-1.5){\pspicture(-1,-1)(1,1)
	
	\pscircle[linecolor=green](0,0){1}
	\pstGeonode[PosAngle={200,-20}](1;-160){A}(1;-20){B}
	\pstGeonode[PosAngle=-90,PointName=C'](1;-90){C}
	
	\pstLineAB{A}{B}
	\pstLineAB{A}{C}
	\pstLineAB{C}{B}
	
	\pstMarkAngle[MarkAngleRadius=0,LabelSep=0.5,LabelAngleOffset=4]{C}{A}{B}{$\alpha'$}
	\pstMarkAngle[MarkAngleRadius=0,LabelSep=0.6]{A}{B}{C}{$\beta'$}
	\pstMarkAngle[MarkAngleRadius=0,LabelSep=0.25]{B}{C}{A}{$\gamma'$}
	\endpspicture
	}
	\pcline[nodesep=1.7\psunit]{->}(0,0)(5,1.5)\aput{:U}{Caso 1}
	\pcline[nodesep=1.7\psunit]{->}(0,0)(5,-1.5)\bput{:U}{Caso 2}
	\endpspicture
\end{center}
Vediamo ora se e quando si pu\`{o} creare un angolo di $60^{\circ}$.
Con la prima trasformazione occorre che sia $\gamma=60^{\circ}$,
oppure $90^{\circ}-\gamma/2=60^{\circ}$, da cui ancora
$\gamma=60^{\circ}$.  Con la seconda trasformazione, occorre che sia
$\gamma=60^{\circ}$ oppure $\gamma=120^{\circ}$. Abbiamo cos\`{\i} 
dimostrato che si pu\`{o} creare un angolo di $60^{\circ}$ se e solo 
se partiamo da una configurazione in cui 2 dei 3 punti sono gi\`{a} 
alla distanza giusta per essere 2 vertici di un triangolo equilatero, 
e cio\`{e} formano angoli alla circonferenza di $60^{\circ}$ e 
$120^{\circ}$, o equivalentemente angoli al centro di $120^{\circ}$ e 
$240^{\circ}$.

Se dunque alla fine abbiamo un triangolo equilatero, procedendo a
ritroso, in ogni configurazione precedente (dunque in particolare
nella configurazione iniziale) dobbiamo avere gi\`{a} 2 punti che
formano angoli al centro di $120^{\circ}$ e $240^{\circ}$.

D'altra parte, se la configurazione iniziale \`{e} di questo tipo, 
possiamo concludere in una sola mossa!}

\renewcommand{\SoluzioniAlternative}{
}

\renewcommand{\Osservazioni}{
}

\renewcommand{\Commenti}{
}